\section{Проведённые эксперименты и~результаты}\label{sec:eval}
\subsection{Условия сравнения}
Для основных сравнений выбран BigCodeBench v0.1.4 в~режиме \texttt{instruct full}~\cite{zhuo2024bigcodebench}. В~каждой задаче модель получает требования к~функции Python и~должна вернуть законченную программу. После получения готовой программы её проверяют исходными тестами набора. Во время работы модель не видит эти тесты, эталонное решение и~результаты итоговой проверки.

Основное сравнение проверяет, что меняется при переходе от последовательных вызовов к~одному вызову. Дополнительные варианты без названий ролей позволяют проверить, связан ли результат с~распределением действий между вызовами или с~ролевыми словами в~запросе. Планирование, написание кода и~рецензирование можно обозначить нейтрально, как этапы 1, 2 и~3, или назвать ролями планировщика, исполнителя и~рецензента. Все три инструкции можно передать одним запросом или распределить между тремя последовательными обращениями. Получаются четыре условия. К~ним добавлен прямой запрос без предписанных этапов, чтобы проверить, нужны ли эти этапы вообще. Все пять условий представлены в~таблице~\ref{tab:design}.

\begin{table}[htbp]
\caption{Условия основного сравнения}
\label{tab:design}
\centering\small\renewcommand{\arraystretch}{1.18}
\begin{tabularx}{\textwidth}{@{}lc>{\raggedright\arraybackslash}p{2.2cm}>{\raggedright\arraybackslash}X@{}}
\toprule
Код & Вызовы & Обозначения & Предписанные действия \\
\midrule
D & 1 & Без этапов & Сразу подготовить программу. \\
SN & 1 & Этапы 1, 2, 3 & Составить план, написать код, проверить и~исправить его по требованиям. \\
SR & 1 & Названия ролей & Планировщик, исполнитель и~рецензент заданы в~одном вызове. \\
MN & 3 & Этапы 1, 2, 3 & Те же действия в~последовательных вызовах. \\
\MultiCallRole{} & 3 & Названия ролей & Планировщик, исполнитель и~рецензент заданы в~трёх отдельных вызовах. \\
\botrule
\end{tabularx}
\end{table}
\FloatBarrier

В~условиях SN, SR, MN и~\MultiCallRole{} формулировки самих действий совпадают. При трёх вызовах каждый получает полный текст задачи и~все предыдущие видимые ответы. Поэтому переход к~трём вызовам одновременно добавляет повторный контекст, промежуточные сообщения и~отдельную генерацию на каждом этапе. Эксперимент оценивает результат такого разделения работы. Определить вклад каждого из этих изменений по отдельности он не позволяет.

Условие SR соответствует интроспективному вызову с~названиями ролей. В~условии \MultiCallRole{} те же три действия с~названиями ролей распределены между последовательными вызовами. Итоговая программа в~обоих случаях проверяется тестами. Сжатие контекста, автоматический выбор числа вызовов и~исправление по результатам внешних тестов в~этом сравнении не используются. Оригинальный INoT также не воспроизводится целиком: исследуется общий для подходов принцип объединения действий внутри одного вызова.

В~первых трёх сериях BigCodeBench использовалась модель с~идентификатором \texttt{gpt-5.4-mini}, в~двух остальных сериях BigCodeBench --- \texttt{gpt-5.6-luna}. Уровень рассуждения во всех пяти сериях BigCodeBench задан как \texttt{medium}. Запросы отправлялись через Codex CLI 0.153.4 с~доступом по подписке. Каждый вызов создавал новую сессию без пользовательских инструкций и~файлов проекта. В~условиях D, SN, SR, MN и~\MultiCallRole{} инструменты были отключены. Внутренние проверки метода SCC описаны в~приложении~\ref{subsec:scc}.

Повторы нужны, потому что одна модель может дать разные программы на один и~тот же запрос. Каждый повтор представлял собой новую генерацию для той же задачи и~условия. Предыдущая программа при этом не исправлялась, а~лучший ответ по скрытым тестам не выбирался. Управление случайным состоянием модели на стороне поставщика не было подтверждено. Общий жёсткий предел выходных токенов не задавался. На один вызов отводилось 180 секунд в~сериях на 40 и~200 задачах и~600 секунд в~сериях на 80 и~1\,000 задачах. Сетевой сбой или превышение времени отмечались как незавершённый запуск. Качество отсутствующей программы оставалось неизвестным.

Основная генерация на 1\,000 задачах проходила 8--11 сентября 2026 года по UTC. Конечные записи выбранных запусков SCC относятся к~11--13 сентября, а~первое сравнение на 200 задачах выполнено 8 сентября. Даты в~именах архивных папок служат идентификаторами архивов и~не всегда совпадают с~датами выполнения. Позднейшие уточнения анализа, включая добавления 3 октября, основаны на этих же сохранённых результатах.

\subsection{Пять серий на BigCodeBench}
Единица плана эксперимента, далее \emph{назначение}, задаёт одну задачу, одно условие и~один повтор. Например, пять условий с~тремя повторами на 1\,000 задачах дают 15\,000 назначений. Таблица~\ref{tab:series} разделяет запланированные генерации, полученные программы и~программы с~доступной оценкой качества.

\begin{table}[htbp]
\caption{Объём пяти серий на BigCodeBench}
\label{tab:series}
\centering\small\renewcommand{\arraystretch}{1.15}
\begin{tabularx}{\textwidth}{@{}>{\raggedright\arraybackslash}Xrrrr@{}}
\toprule
Серия & Задачи & Назначено & Получено & Оценено \\
\midrule
Отладочная & 40 & 400 & 400 & 390 \\
Первое сравнение & 200 & 1\,000 & 996 & 963 \\
Проверка оформления & 80 & 320 & 320 & 316 \\
Расширенное сравнение & 1\,000 & 15\,000 & 14\,961 & 14\,736 \\
Сравнение с~SCC & 956 & 6\,000 & 5\,595 & 5\,503 \\
\botrule
\end{tabularx}
\end{table}
\FloatBarrier

На 40 задачах отладили весь цикл: извлекали из ответа модели программу на Python, запускали её в~изолированной среде на тестах BigCodeBench и~записывали, прошла ли она тесты. В~серии было пять условий по два повтора. Её результаты приведены отдельно от основных сравнений токенного расхода и~качества. Первое основное сравнение проводилось на 200 других задачах: пять условий по одной новой генерации. Задачи выбирались из случайной перестановки допустимых идентификаторов. Начальное состояние генератора Python было равно 20260908. Если контроль среды выявлял проблему, задача сохранялась в~плане и~не заменялась более удобной.

На 80 новых задачах сравнили два способа записи инструкций внутри одного вызова. В~первом каждое действие начиналось с~новой строки и~метки в~квадратных скобках (приложение~\ref{app:prompt-lines}). Во втором те же действия шли подряд в~одном абзаце, с~двоеточием после каждой метки (приложение~\ref{app:prompt-paragraph}). В~каждом формате были две версии: с~метками \texttt{PLANNER}, \texttt{IMPLEMENTER}, \texttt{REVIEWER} и~с~нейтральными метками \texttt{STAGE 1}, \texttt{STAGE 2}, \texttt{STAGE 3}. Обе версии поручали модели составить план, написать функцию и~проверить собственный код. Внешние тесты оценивали готовую программу, поэтому по результатам можно сравнивать формулировки запроса, а~не работу каждого этапа по отдельности. Четыре сочетания дали 320 программ. Эти 80 задач не пересекаются с~отладочными 40 и~первыми 200.

Расширенная серия включала 1\,000 задач: прежние 200 и~ещё 800. Для каждой задачи и~каждого из пяти условий заранее назначено по три новые генерации. В~протоколе это число выбрано, чтобы сохранить охват 1\,000 задач в~пределах бюджета. Вариант с~пятью повторами требовал бы 25\,000 назначений вместо 15\,000. Оптимальность именно трёх повторов не установлена. Это число не гарантирует заданной точности. Их результаты усредняются внутри задачи, а~полученная точность оценивается доверительными интервалами. Старые программы не использовались. Таким образом, увеличились число задач и~число повторов, а~модель mini была заменена на Luna. Сопоставлять серии как проверку только одного из этих изменений нельзя.

Для сравнения с~SCC выбраны первые 2\,000 сочетаний «задача--повтор» из случайно перемешанного порядка. Каждому сочетанию назначены SN, SR и~SCC, всего 6\,000 запусков. Эти сочетания охватывают 956 задач из расширенной серии: у~212 задач выбран один повтор, у~444 --- два, у~300 --- три. Исходный план на 9\,000 назначений был сокращён после начала генерации, но до просмотра качества SCC. План и~поправка сохранены. Во всех пяти сериях вместе встречается 1\,120 разных идентификаторов BigCodeBench. Суммировать размеры пересекающихся выборок было бы неверно.

\subsection{Проверка программ и~работа с~пропусками}\label{subsec:missing-outcomes}
Прежде чем оценивать ответы модели, для каждой задачи проверяли эталонную и~заведомо неверную программы. Задача включалась в~анализ качества, если эталон проходил тесты, а~неверная программа отклонялась. Такой контроль помогает отделить ошибку ответа от проблемы окружения. Программы исполнялись штатным оценщиком BigCodeBench в~изолированном контейнере без сети. Исходные тесты не менялись под ответы модели. Для каждой программы сохранялись результат, сообщения проверяющего и~связь с~назначением.

Доля успехов равна числу программ, прошедших все тесты, делённому на число программ с~известной оценкой. Оценка могла отсутствовать по двум причинам: генерация не завершилась либо среда не прошла контроль на этой задаче. Эти случаи записывались раздельно и~исключались из расчёта доли успехов по известным исходам. Программы, которые удалось проверить и~которые не прошли тесты, оставались в~анализе как неуспешные. Поэтому исключение неизвестных результатов не означает исключения ошибочного кода. Известные расходы незавершённых запусков сохранялись. Отсутствующий счётчик токенов отмечался как неизвестный.

Сравнение двух методов строилось на общих задачах. В~расширенной серии для основной оценки требовались все три результата в~обоих условиях. Сначала три повтора усреднялись внутри задачи, затем вычислялась разность между условиями. Итог представлял собой среднее этих разностей с~одинаковым весом каждой задачи. Поэтому три повтора не считались тремя независимыми задачами, а~число доступных пар могло различаться между сравнениями.

Четыре сравнения качества заданы заранее: SR с~SN, \MultiCallRole{} с~MN, MN с~SN и~\MultiCallRole{} с~SR. В~серии на 200 задачах применён точный двусторонний критерий Мак-Немара. В~серии на 1\,000 задачах использованы двусторонние парные $t$-тесты для средних по задачам. Решения принимались с~поправкой Холма~\cite{holm1979}. При предпосылках тестов она ограничивает вероятность хотя бы одного ложного положительного вывода в~этом семействе величиной 0,05.

Рядом с~оценками приведены описательные 95\%-е интервалы. Для их расчёта 10\,000 раз выбирались с~возвращением целые задачи со всеми повторами. Такой бутстрэп показывает неопределённость оценки, сохраняя связь повторов одной задачи. Эти интервалы не учитывают четыре проверки одновременно. Решение о~значимости принимается по тестам с~поправкой Холма.

Основная оценка использует задачи с~известными исходами. Чтобы проверить устойчивость вывода на всей назначенной выборке, дополнительно находим крайние возможные значения средней разности. Пусть $D$ --- разность числа известных успехов способов $A$ и~$B$, $U_A$ и~$U_B$ --- числа прогонов без оценки, а~$N$ --- число назначенных прогонов на способ. Нижняя граница равна $100(D-U_B)/N$ п.п. Она получается, если все неизвестные исходы $B$ окажутся успехами, а~все неизвестные исходы $A$ --- неуспехами. При обратном распределении получается верхняя граница $100(D+U_A)/N$. Минимум и~максимум наблюдавшихся исходов характеризуют разброс отдельных наблюдений. По рассчитанным границам видно, насколько пропуски могут сдвинуть среднюю разность на всех назначенных задачах.

Исходный план расширенной серии был сохранён до начала её генерации, а~контрольные проверки проведены до просмотра качества ответов. Это локальная фиксация протокола. Внешняя предварительная регистрация исследования не проводилась. После прерываний разрешалось продолжать только назначения, в~которых ещё не начался ни один вызов. Начатый неудачный запуск не заменялся новой попыткой в~той же ячейке. Порядок назначений перемешивался, но фактическое время завершения зависело также от параллельного выполнения и~сбоев.

\subsection{Качество на 1\,000 задачах}\label{subsec:scale}
Из 15\,000 назначений завершилось 14\,961. 39 остались незавершёнными. Контроль среды позволил оценивать качество на 985 задачах. В~результате известны 14\,736 исходов, а~264 остаются неизвестными из-за незавершённой генерации или непригодности среды. Таблица~\ref{tab:scale-quality} показывает результат каждого условия. Последний столбец рассчитан от числа оценённых программ в~той же строке.

\begin{table}[htbp]
\caption{Качество в~расширенной серии: исходы отдельных генераций}
\label{tab:scale-quality}
\centering\small\renewcommand{\arraystretch}{1.15}
\begin{tabular}{@{}lrrrr@{}}
\toprule
Условие & Получено & Оценено & Успехов & Успех, \% \\
\midrule
D & 2\,996 & 2\,951 & 1\,611 & 54,59 \\
SN & 2\,998 & 2\,953 & 1\,604 & 54,32 \\
SR & 2\,997 & 2\,952 & 1\,581 & 53,56 \\
MN & 2\,983 & 2\,938 & 1\,471 & 50,07 \\
\MultiCallRole{} & 2\,987 & 2\,942 & 1\,514 & 51,46 \\
\botrule
\end{tabular}
\end{table}
\FloatBarrier

Каждому условию было назначено 3\,000 генераций. Доступные исходы в~строках различаются, поэтому простое вычитание процентов из таблицы~\ref{tab:scale-quality} не даёт парной оценки. В~таблице~\ref{tab:scale-contrasts} разности рассчитаны на одних и~тех же задачах с~полными повторами. В~последних двух строках показано «один вызов минус три»: положительная разность означает преимущество одного вызова. Перемена знака относительно исходной записи сравнения не меняет двусторонний тест. Разность выражена в~процентных пунктах, сокращённо п.п.

\begin{table}[htbp]
\caption{Парные различия качества на 1\,000 задачах}
\label{tab:scale-contrasts}
\centering\small\renewcommand{\arraystretch}{1.2}
\begin{tabular}{@{}lrrrr@{}}
\toprule
Сравнение & Задачи & Разность, п.п. & 95\% интервал & $p$ Холма \\
\midrule
SR $-$ SN & 982 & $-0{,}75$ & $[-1{,}97;+0{,}44]$ & 0,23575 \\
\MultiCallRole{} $-$ MN & 964 & $+1{,}35$ & $[+0{,}17;+2{,}56]$ & 0,05986 \\
SN $-$ MN & 968 & $+4{,}17$ & $[+2{,}75;+5{,}65]$ & $1{,}34\cdot10^{-7}$ \\
SR $-$ \MultiCallRole{} & 971 & $+2{,}13$ & $[+0{,}82;+3{,}43]$ & 0,00445 \\
\botrule
\end{tabular}
\end{table}
\FloatBarrier

Доля программ, прошедших тесты, у~интроспективного вызова SR была на 2,13 процентного пункта выше, чем у~трёх последовательных вызовов \MultiCallRole{}. Для нейтральных инструкций разность между SN и~MN составила 4,17 пункта в~пользу SN. Оба результата проходят поправку Холма. Знак сохраняется даже при заполнении всех пропусков в~пользу трёх вызовов. По всей тысяче задач крайние границы равны $[+0{,}30;+3{,}83]$ пункта для SR минус \MultiCallRole{} и~$[+2{,}37;+6{,}00]$ для SN минус MN.

Сами названия ролей не дали подтверждённого улучшения. Различие SR и~SN статистически незначимо. Для \MultiCallRole{} минус MN оценка положительна, но скорректированное $p=0{,}05986$ превышает 0,05. Здесь описательный интервал не пересекает ноль, однако он не учитывает одновременную проверку четырёх гипотез. Этот пример показывает, зачем различать интервал отдельной оценки и~решение с~поправкой. При этом отсутствие значимости оставляет возможность небольшого эффекта.

\subsection{Проверки устойчивости основного результата}
После основного анализа выполнено несколько дополнительных проверок на сохранённых данных. Они помогают понять границы вывода, но не меняют четыре заранее заданных теста и~не добавляют новых генераций.

Сначала для тех же четырёх разностей построены более строгие интервалы, учитывающие их совместную оценку. Каждый интервал Стьюдента имеет уровень 98,75\%. Правило Бонферрони даёт приблизительное совместное покрытие не менее 95\% при предпосылках анализа по задачам. Интервалы показаны в~таблице~\ref{tab:precision}. Оба различия между одним и~тремя вызовами остаются положительными. Для названий ролей интервалы включают ноль и~допускают небольшой выигрыш: до 0,83 пункта при одном вызове и~до 2,90 при трёх.

\begin{table}[htbp]
\caption{Дополнительные интервалы с~учётом четырёх сравнений}
\label{tab:precision}
\centering\small
\begin{tabular}{@{}lrr@{}}
\toprule
Сравнение & Задачи & 98,75\% интервал, п.п. \\
\midrule
SR $-$ SN & 982 & $[-2{,}32;+0{,}83]$ \\
\MultiCallRole{} $-$ MN & 964 & $[-0{,}20;+2{,}90]$ \\
SN $-$ MN & 968 & $[+2{,}29;+6{,}04]$ \\
SR $-$ \MultiCallRole{} & 971 & $[+0{,}46;+3{,}80]$ \\
\botrule
\end{tabular}
\end{table}
\FloatBarrier

Затем результаты разделены по времени генерации. Первый завершённый сеанс содержит 5\,551 программу, последующие --- 9\,410. В~каждом периоде сопоставлялись только повторы с~известными результатами обоих методов. Разности усреднялись сначала внутри задачи, затем по задачам. Таблица~\ref{tab:temporal} показывает все четыре сравнения. Оценки преимущества одного вызова положительны в~обоих периодах. Однако ранний интервал SR минус \MultiCallRole{} включает ноль, то есть данных этого среза недостаточно для уверенного вывода о~знаке различия. Периоды пересекаются по задачам и~не являются независимыми повторениями эксперимента. Такая проверка также не исключает изменений модели у~поставщика.

\begin{table}[htbp]
\caption{Различия качества в~разные периоды генерации}
\label{tab:temporal}
\centering\small\setlength{\tabcolsep}{4pt}
\begin{tabular}{@{}llrrl@{}}
\toprule
Сравнение & Период & Задачи & \shortstack{Пары\\повторов$^*$} & \shortstack{Разность, п.п.\\{}[95\% интервал]} \\
\midrule
SR $-$ SN & Первый & 740 & 1\,096 & $-1{,}42\;[-3{,}54;+0{,}74]$ \\
 & Поздние & 936 & 1\,855 & $-0{,}66\;[-2{,}23;+0{,}91]$ \\
\MultiCallRole{} $-$ MN & Первый & 736 & 1\,090 & $+0{,}79\;[-1{,}40;+2{,}99]$ \\
 & Поздние & 934 & 1\,844 & $+1{,}41\;[-0{,}11;+2{,}93]$ \\
SN $-$ MN & Первый & 737 & 1\,092 & $+4{,}27\;[+1{,}99;+6{,}56]$ \\
 & Поздние & 934 & 1\,846 & $+4{,}46\;[+2{,}61;+6{,}37]$ \\
SR $-$ \MultiCallRole{} & Первый & 738 & 1\,092 & $+2{,}06\;[-0{,}14;+4{,}27]$ \\
 & Поздние & 934 & 1\,849 & $+2{,}37\;[+0{,}77;+4{,}00]$ \\
\botrule
\end{tabular}
\par\smallskip\raggedright\footnotesize $^*$Число пар повторов с~результатами обоих методов. Интервалы описательные, получены 10\,000 выборками целых задач, PCG64, начальное состояние 20261003.
\end{table}
\FloatBarrier

Ещё одна проверка учитывает историю выборки. Из 1\,000 задач 200 уже встречались в~ранней серии, а~800 использовались впервые. В~таблице~\ref{tab:subgroups} эти части рассмотрены отдельно. Преимущество одного вызова сохраняется на 800 новых идентификаторах: 4,17 пункта при нейтральных инструкциях и~2,35 при ролевых. Следовательно, основной результат не определяется только прежними 200 задачами. Все программы, включая ответы на прежние задачи, были сгенерированы заново.

\begin{table}[htbp]
\caption{Различия качества на прежних и~новых задачах}
\label{tab:subgroups}
\centering\small\setlength{\tabcolsep}{5pt}
\begin{tabular}{@{}llrl@{}}
\toprule
Сравнение & Выборка & Задачи & \shortstack{Разность, п.п.\\{}[95\% интервал]} \\
\midrule
SR $-$ SN & Прежние 200 & 194 & $-1{,}03\;[-3{,}78;+1{,}72]$ \\
 & Новые 800 & 788 & $-0{,}68\;[-2{,}07;+0{,}72]$ \\
\MultiCallRole{} $-$ MN & Прежние 200 & 190 & $+1{,}93\;[-0{,}70;+4{,}56]$ \\
 & Новые 800 & 774 & $+1{,}21\;[-0{,}13;+2{,}58]$ \\
SN $-$ MN & Прежние 200 & 192 & $+4{,}17\;[+0{,}87;+7{,}64]$ \\
 & Новые 800 & 776 & $+4{,}17\;[+2{,}58;+5{,}84]$ \\
SR $-$ \MultiCallRole{} & Прежние 200 & 190 & $+1{,}23\;[-1{,}23;+3{,}68]$ \\
 & Новые 800 & 781 & $+2{,}35\;[+0{,}81;+3{,}88]$ \\
\botrule
\end{tabular}
\end{table}
\FloatBarrier

Наконец, проверено сходство самих задач. Для инструкций сравнивались множества последовательностей из пяти слов после удаления служебного заключительного абзаца и~приведения к~одному регистру. Мера Жаккара равна числу общих элементов, делённому на число разных элементов в~объединении двух множеств. Инструкции связывались при значении не ниже 0,70. Отдельно проверялись пороги 0,50 и~0,90.

Объединённое правило также учитывало точные совпадения инструкций после нормализации пробелов, совпадающую структуру эталонного кода без строк документации и~близкие наборы имён и~литералов кода. Для последнего признака требовалось не менее 20 вхождений идентификаторов в~каждой программе, мера Жаккара не ниже 0,80 и~её вариант с~учётом повторений не ниже 0,70. Связи строились по всем 1\,140 задачам BigCodeBench. Задачи, связанные напрямую или через промежуточные задачи, объединялись в~группы. После ограничения на выбранную тысячу получилось 992 группы.

При повторной выборке целых групп с~сохранением весов задач интервалы для SR минус SN и~\MultiCallRole{} минус MN составили $[-1{,}99;+0{,}47]$ и~$[+0{,}14;+2{,}60]$ пункта. Для SN минус MN и~SR минус \MultiCallRole{} они равны $[+2{,}72;+5{,}69]$ и~$[+0{,}79;+3{,}44]$. Использованы 10\,000 выборок, PCG64 и~начальное состояние 20260911. Направление основного результата сохранилось при таком учёте сходства. Протокол этой проверки записан 11 сентября, уже во время генерации. Проверка охватывает заданные виды сходства. Она не устанавливает независимость всех задач или отсутствие этих задач в~обучающих данных модели.

\subsection{Сравнение с~прямым запросом}
Прямой запрос D даёт ориентир для оценки всех дополнительных этапов. Его парные сравнения в~таблице~\ref{tab:direct} выполнены после основного анализа и~носят разведочный характер. Они не входят в~семейство четырёх тестов с~поправкой Холма.

\begin{table}[htbp]
\caption{Разведочные сравнения качества прямого запроса}
\label{tab:direct}
\centering\small
\begin{tabular}{@{}lrrr@{}}
\toprule
Сравнение & Задачи & Разность, п.п. & 95\% интервал \\
\midrule
D $-$ SN & 980 & $+0{,}27$ & $[-0{,}88;+1{,}43]$ \\
D $-$ SR & 979 & $+0{,}95$ & $[-0{,}17;+2{,}11]$ \\
D $-$ MN & 966 & $+4{,}59$ & $[+2{,}97;+6{,}18]$ \\
D $-$ \MultiCallRole{} & 970 & $+3{,}20$ & $[+1{,}82;+4{,}64]$ \\
\botrule
\end{tabular}
\end{table}
\FloatBarrier

Различия D с~одним нейтральным или ролевым вызовом невелики, а~оба интервала включают ноль. Это оставляет неопределённость относительно их порядка и~не доказывает равенства. По сравнению с~тремя вызовами наблюдаемая доля успехов D выше на 3,20--4,59 пункта. Поэтому объединение этапов полезно оценивать и~относительно простого прямого запроса: SR нельзя считать лучшим из всех пяти условий только по его преимуществу над \MultiCallRole{}.

В~сравнении с~многошаговой процедурой SCC разность долей успеха SCC и~SR на 883 задачах с~общими доступными повторами составила $-2{,}68$ процентного пункта. Интервал, учитывающий все неизвестные исходы на 941 пригодной задаче, допускает оба направления эффекта. На 896 задачах с~полными парными счётчиками расчётная стоимость SCC была в~3,14 раза выше. Реализация и~полный анализ приведены в~приложении~\ref{subsec:scc}.

\subsection{Расход токенов и~расчётная стоимость}
В~расход включены входы и~выходы всех вызовов, в~том числе повторный контекст и~промежуточные ответы. Кэшированные токены уже входят в~общий вход, а~токены рассуждения --- в~выход. Повторно они не прибавляются.

Денежная оценка переводит эти счётчики в~стоимость по сохранённым тарифам программного интерфейса сервиса, или API. За миллион обычных входных, кэшированных входных и~выходных токенов использованы соответственно 0,75 / 0,075 / 4,50 доллара для mini и~0,20 / 0,02 / 1,20 доллара для Luna. Это фиксированные коэффициенты расчёта. Поскольку эксперимент выполнен через подписку, полученные суммы не являются фактическими платежами за API или стоимостью подписки. Время исследователя и~вычисления оценщика в~них не включены.

Ресурсы сравниваются на задачах со всеми тремя повторами и~известными счётчиками обоих методов. Сначала вычисляются средние внутри задач, затем средние по общему набору задач. Их отношение приведено в~таблице~\ref{tab:scale-resources}. Например, значение 2,780 для \MultiCallRole{}/SR означает, что три последовательных вызова имеют в~2,780 раза большую тарифную оценку, чем один интроспективный вызов. Наборы полных пар для качества и~ресурсов могут различаться.

\begin{table}[htbp]
\caption{Отношения ресурсов: три вызова к~одному}
\label{tab:scale-resources}
\centering\small\renewcommand{\arraystretch}{1.2}
\begin{tabular}{@{}lrrrr@{}}
\toprule
Сравнение & Задачи & Токены & API-оценка & Без скидки за кэш \\
\midrule
MN / SN & 983 & 3,078 & 2,672 & 2,777 \\
\MultiCallRole{} / SR & 986 & 3,090 & 2,780 & 2,867 \\
\botrule
\end{tabular}
\end{table}
\FloatBarrier

Три последовательных вызова \MultiCallRole{} потребовали в~3,09 раза больше токенов и~в~2,78 раза больших расчётных затрат, чем один интроспективный вызов SR. В~обратном сравнении SR использовал \textbf{на 67,6\% меньше токенов}, чем \MultiCallRole{}. Описательный 95\%-й интервал отношения затрат \MultiCallRole{}/SR равен $[2{,}74;2{,}82]$, а~MN/SN --- $[2{,}63;2{,}71]$. Преимущество одного вызова по расчётным затратам сохраняется и~при одинаковой цене всего входа, без скидки на кэш.

Абсолютная разница помогает оценить практический масштаб этих отношений. На полных ресурсных парах три нейтральных вызова добавляют в~среднем 8\,582 токена на программу по сравнению с~одним, а~три ролевых --- 8\,510. В~обоих случаях разница тарифных оценок составляет \textbf{примерно 1,60 доллара на 1\,000 генераций}. Для этих коротких задач большое относительное различие соответствует небольшой денежной сумме по использованным тарифам. Время измерялось отдельно в~локальном пилоте (раздел~\ref{subsec:route-pilot}).

Известная тарифная оценка всех завершённых кандидатов пяти условий основной серии составляет 23,07 доллара. Известные этапы незавершённых запусков добавляют около 0,03 доллара, всего 23,10 доллара. Счётчики 32 начатых этапов недоступны и~в~сумму не входят. Эти общие суммы показывают масштаб исследования. Сравнение условий опирается на полные пары, чтобы сопоставлять расходы на одном и~том же наборе задач. Такое ограничение само по себе не устраняет возможное смещение из-за пропусков.

Названия ролей меняют ресурсы заметно слабее. Отношение тарифных оценок SR/SN равно 0,939 на 997 задачах, \MultiCallRole{}/MN --- 0,978 на 979. Это описательные различия, которые нужно рассматривать вместе с~неопределённостью качества. Длина контекста отдельно не изменялась, поэтому порог из утверждения~\ref{prop:savings} по этой серии не оценивался.

\subsection{Сопоставление качества и~расхода на общих задачах}
Из 986 задач с~полными данными о~расходе 15 не прошли контроль тестовой среды. Для сравнения качества и~расходов вместе используем 971 задачу, для которой известны исходы тестов, счётчики токенов и~расчётная стоимость всех трёх прогонов каждого способа. В~общий расход входят все прогоны из этой выборки, включая программы, провалившие тесты. Таблица~\ref{tab:scale-quality-cost} показывает оба показателя по одному и~тому же набору задач.

\begin{table}[htbp]
\caption{Качество и~расходы на общей выборке из 971 задачи}
\label{tab:scale-quality-cost}
\centering\small\renewcommand{\arraystretch}{1.18}
\begin{tabular}{@{}lrr@{}}
\toprule
Показатель & SR & \MultiCallRole{} \\
\midrule
Число прогонов & 2\,913 & 2\,913 \\
Число программ, прошедших тесты & 1\,562 & 1\,500 \\
Доля успешных программ, \% & 53,62 & 51,49 \\
Токенов на прогон, в~среднем & 4\,070 & 12\,578 \\
Расчётная стоимость всех прогонов, долл. & 2,60695 & 7,25753 \\
Расчётная стоимость на успешную программу, долл. & 0,001669 & 0,004838 \\
\botrule
\end{tabular}
\end{table}
\FloatBarrier

Последняя строка --- суммарная расчётная стоимость, делённая на число успешных программ, или $1/U_{\$}$ из формулы~\eqref{eq:efficiency}. Если средняя стоимость прогона и~доля успехов уменьшатся в~одинаковой пропорции, стоимость на успешную программу останется прежней. В~этой выборке SR повысил долю успехов на 2,13 процентного пункта, а~\textbf{стоимость на успешную программу снизилась на 65,5\%}. Это описательный пересчёт на задачах с~полными данными. Тарифные оценки 15 задач с~непроверяемым качеством учтены при оценке расходов на 986 задачах.

\subsection{Результаты первых трёх серий}
В~отладочной серии контроль среды прошли 39 задач. Для каждого условия доступны 78 исходов, из которых успешны 42 для D, 34 для SN, 39 для SR, 39 для MN и~36 для \MultiCallRole{}. Эти результаты использовались при разработке процедуры и~не объединяются с~последующими сериями для проверки гипотез.

В~первом основном сравнении на 200 задачах доли успехов составили 94/193 для D, 93/191 для SN, 92/193 для SR, 86/193 для MN и~88/193 для \MultiCallRole{}. Здесь числитель показывает успешные программы, а~знаменатель --- доступные оценки. Четыре назначения не завершились. Семь задач не прошли контроль среды. Ни одно из четырёх основных сравнений качества не прошло поправку Холма. Например, SR минус SN равно $-1{,}05$ пункта с~интервалом $[-4{,}71;+2{,}62]$. Тарифная оценка SR на парной выборке была на 18,45\% ниже, чем SN, однако разность качества оставалась неопределённой.

В~серии на 80 задачах сравнили метки ролей и~номера этапов отдельно для каждого формата при одних и~тех же требованиях к~действиям модели. Из 79 пригодных задач при записи инструкций на отдельных строках оба варианта решили по 44. Разность равна нулю, 95\%-й интервал составляет $[-6{,}33;+6{,}33]$ пункта. При записи одним абзацем вариант с~названиями ролей решил 42 задачи, с~номерами этапов --- 43. Разность равна $-1{,}27$ пункта с~интервалом $[-10{,}13;+6{,}33]$. Оба значения $p$ после поправки Холма равны 1. Вклад названий ролей при неизменных этапах остаётся неопределённым. Интервалы допускают и~улучшение, и~снижение доли успешных решений на несколько процентных пунктов. 

\subsection{Гибридный выбор числа вызовов}\label{subsec:route-pilot}
Для RQ4 провели отдельную серию с~локальной Llama~3.2~3B Instruct (квантование Q4\_K\_M, Ollama 0.40.16, NVIDIA GeForce GTX~1060 с~6~ГБ памяти). Для неё составили 12 новых коротких задач на целочисленные функции Python. Например, требовалось вычислить смещение страницы, определить длину пересечения интервалов или округлить длительность в~секундах до целых минут в~сторону нуля. В~архивных карточках 1\,000 задач BigCodeBench отсутствовало поле с~требованием к~сессии исполнителя. В~новой серии это требование задали явно.

До генерации для шести задач в~протоколе задали требование отдельной сессии рецензента. Для остальных шести разрешили объединить этапы. Это организационное условие считывается из карточки задачи. Оно задаёт состав нагрузки пилота и~не оценивает сложность задачи. Требования параллельного выполнения в~наборе отсутствовали. Гибрид выбирал один или три вызова. Два контрольных способа всегда делали один или всегда три. Все три способа использовали одни и~те же названия ролей. Ещё в~двух условиях названия ролей заменили номерами этапов, сохранив формулировки действий и~число вызовов.

Каждая задача прошла все пять условий по одному разу. Порядок задач перемешали. Порядок условий внутри задачи сдвигали на одну позицию от задачи к~задаче, начиная со~случайной перестановки. Начальное состояние генератора зафиксировали, температуру задали равной нулю. В~отдельной сессии модель получала условие и~явно переданные ответы предшествующих этапов. На отдельный план отводилось до 40 выходных токенов, на черновой или итоговый код --- до~112. Гибрид для каждого назначения заново обращался к~модели. Всего выполнено 60 прогонов и~120 измеренных обращений после одного прогрева.

Во всех условиях модель должна была вернуть функцию Python с~одним оператором \texttt{return}. В~инструкции перечислялись допустимые операции над числами. Независимый оценщик проверял форму ответа и~совпадение результата с~ожидаемым целым числом на всех входах задачи. Всего в~12 наборах было 648 контрольных входов, включая граничные случаи. Тесты оставались скрытыми от модели. Время измеряли от выбора способа до окончания проверки, включая промежуточные запросы. Все обращения выполнялись последовательно на прогретой модели. Кэш между условиями сохранялся.

\begin{table}[!ht]
\caption{Локальный пилот. Показатели рассчитаны по одним и~тем же 12 задачам в~каждом условии}\label{tab:pilot}
\centering\small\renewcommand{\arraystretch}{1.22}
\begin{tabular}{@{}lrrrrr@{}}
\toprule
Способ & Успехи & Вызовы & Токены & Долл./1\,000 & Время, с \\
\midrule
Один с~ролями (SR) & 5/12 & 12 & 186,8 & 0,0612 & 0,66 \\
Гибрид (RouteMode) & 5/12 & 24 & 384,9 & 0,1272 & 1,35 \\
Три с~ролями (MAS) & 5/12 & 36 & 572,9 & 0,1876 & 1,93 \\
\midrule
Один с~номерами (SN) & 4/12 & 12 & 192,9 & 0,0636 & 0,68 \\
Три с~номерами (MN) & 4/12 & 36 & 608,2 & 0,2048 & 2,20 \\
\botrule
\end{tabular}
\end{table}
\FloatBarrier

Столбец <<Вызовы>> содержит общее число обращений за 12 прогонов. Токены (вход и~выход по счётчикам модели) и~время усреднены на один прогон, включая неуспешные. Для сопоставления цены взяты одинаковые условные ставки из расчёта без скидки на кэш, 0,20~доллара за миллион входных токенов и~1,20~доллара за миллион выходных. Стоимость дана в~пересчёте на 1\,000 прогонов. Генерация выполнена локально, денежные суммы представляют тарифный сценарий.

\textbf{Гибрид и~оба фиксированных способа с~ролями прошли тесты на одних и~тех же пяти задачах.} По сравнению с~постоянными тремя вызовами гибрид использовал в~среднем на 188 токенов меньше на задачу и~завершался на 0,58~с быстрее. По сравнению с~постоянным одним вызовом он израсходовал на 198 токенов больше и~занял на 0,70~с больше. Эти величины относятся к~нагрузке, в~которой отдельная сессия требовалась ровно в~половине случаев. Совпадение исходов на 12 задачах с~единственной генерацией на условие не устанавливает равенства качества на других задачах. Абсолютное время зависит от оборудования и~состояния кэша.

В~дополнительном сравнении названий ролей с~номерами этапов результат изменился на одной задаче из 12 в~обоих сравнениях. С~названиями ролей программа корректно вычислила номер позиции на кольце. При трёх нейтральных вызовах ответ дал неверное значение при положительном числе позиций в~кольце. При одном нейтральном вызове программа совпала с~эталоном на всех входах, но использовала два оператора \texttt{return}. Дополнительная проверка с~разрешённым ранним возвратом повысила результат SN с~4/12 до~5/12, показатели остальных условий сохранились. Таким образом, для одного вызова наблюдаемое различие относилось к~соблюдению формата ответа. На остальных 11 задачах исходы совпали. В~серии BigCodeBench соответствующие эффекты названий ролей составили $-0{,}75$ и~$+1{,}35$ процентного пункта (таблица~\ref{tab:scale-contrasts}). Это разные оценки для разных моделей и~наборов задач. Данные, полные запросы и~программа пересчёта приведены в~приложении~\ref{app:materials}.

\section{Область применимости и~ограничения}\label{sec:limits}
Расчётная стоимость охватывает обращения к~модели. Состав полной стоимости эксплуатации, включая оборудование, электроэнергию и~сопровождение, разобран в~приложении~\ref{app:tco}.

Результаты основной серии относятся к~созданию функций Python из BigCodeBench при выбранных моделях и~настройках. Исправление большого репозитория требует поиска файлов, понимания связей между ними и~проверки изменений. Полные репозитории в~выборку не входили. Публичность BigCodeBench также оставляет открытым вопрос о~присутствии его задач в~обучающих данных.

При разделении на три вызова менялись передача истории и~повторный контекст. Вместо одной генерации выполнялись три. Сравнение поэтому оценивает конкретный способ организации работы. В~нём не проверены равный вычислительный бюджет, отдельный вклад каждого механизма и~перенос на независимые семейства моделей. Дополнительный вызов с~новыми сведениями, инструментом или полезной обратной связью может дать иной результат.

Название модели в~сервисе не фиксирует её неизменяемую версию. Время запуска, состояние кэша и~изменения инфраструктуры могут влиять на наблюдения. Сравнение периодов показывает, сохранялось ли направление различий в~разное время генерации. Оно не позволяет установить, менял ли поставщик модель. Для SCC особенно существенны пропуски из-за контейнеров: их учёт меняет допустимый вывод о~качестве всей выборки.

Наконец, самопроверка может повторять ошибку первоначального ответа, а~сжатие --- удалить нужную деталь. По этой причине Hybrid-INoT предусматривает внешнюю проверку как критерий приёмки программы. Даже её положительный результат относится только к~выполненным тестам и~не доказывает правильность на всех входах. Сгенерированный код исполняется в~ограниченной среде. Ролевые инструкции сами по себе такой защиты не создают.

\section{Заключение}
В~серии BigCodeBench на 1\,000 задачах планирование, написание кода и~рецензирование за один интроспективный вызов потребовали на 67,6\% меньше токенов, чем те же действия за три последовательных вызова с~полной передачей контекста. На 986 задачах с~полными счётчиками три вызова обошлись в~2,78 раза дороже по тарифной оценке. На общей выборке из 971 задачи независимые тесты прошли 1\,562 из 2\,913 программ при одном вызове и~1\,500 из 2\,913 при трёх. Разность в~2,13 процентного пункта и~экономия в~65,5\% на одну успешно проверенную программу характеризуют конкретный способ распределения работы между запросами.

Направление сравнения сохранялось и~при нейтральных названиях этапов. Преимущество одного вызова над тремя по доле успешных программ составило 4,17 процентного пункта с~номерами этапов и~2,13 пункта с~названиями ролей. Таким образом, преимущество объединения этапов по доле успехов наблюдался при обоих способах записи инструкции.

В~локальном пилоте гибридный контроллер направил шесть задач в~отдельные сессии и~шесть в~один вызов. Все три ролевых способа решили одни и~те же 5 из~12 задач. Среднее время составило 0,66~с для одного вызова, 1,35~с для гибрида и~1,93~с для трёх вызовов. В~этих условиях объединение этапов было самым экономным, а~выбор отдельной сессии только по требованию задачи позволил сократить средний расход по сравнению с~назначением трёх вызовов всем задачам.
